\scanpage{310}\mspage{34}
flèches les seuls foncteurs cartésiens. (On fera
attention cependant que le foncteur canonique
$\mathrm{Fibsc}\,S\mathrel{\unsure{\text{---}|\longrightarrow}}\mathrm{Fib}\,S$ est fidèle mais non pleinement ; néanmoins il induit une équivalence sur les catégories localisées, \deleted{\unsure{)}} en
définissant de façon évidente, par \unsure{la|les} $W$-\unsure{équivalence|équivalences} fibre par fibre, i.e comme $W_S^u$,
le localiseur \inserted{pertinent} sur $\mathrm{Fib}\,S$.) Donc pour cette
raison de principe, et \deleted{\unsure{d'a}}d'ailleurs tout au
long du développement de la théorie
de $\mathrm{HOT}_W$, la catégorie $\mathrm{Fib}\,S$ et par là,
les foncteurs fibrants, ont joué un rôle
essentiel. La \deleted{catégorie} \inserted{classe de flèches} que nous notons ici
$\deleted{\unsure{\mathcal F}}\mathcal F_W$ ou $\mathrm{Catfib}^{W}_{\deleted{\unsure{W}}}$ (on pourrait aussi la
noter $\mathrm{Fib}^{W}_{\deleted{\unsure{W}}}$, \deleted{\unsure{si}} ce n'était pour les risques
de confusion avec $\mathrm{Fib}_W$) \deleted{\unsure{est}} s'\deleted{\unsure{ex}}\unsure{interprète|explicite}
alors très naturellement à partir de là.

Une flèche $f:X\longrightarrow S$ est dans $\mathcal F_W$ sss
elle est \quad a) fibrante, et par là définit
de façon particulièrement directe un objet
de $\mathrm{HOT}_W(S)$, puisque à $S$-équivalence près, $X/S$
provient d'un \deleted{\unsure{élément}} \inserted{objet} de $\mathrm{Fibsc}\,S$ ou encore de $\mathrm{Cat}(S)$),
% archive sheet number: 309
